\documentclass[pdf,notes]{beamer}

\mode<presentation>{\usetheme[secheader]{Boadilla}}
%%\setbeamertemplate{theorems}[numbered]
%%\newtheorem{thm}{Theorem}
%%\newtheorem{proposition}[thm]{Proposition}
%%\theoremstyle{example}
%%\newtheorem{exam}[thm]{Example}
%%\newtheorem{remark}[thm]{Remark}
%%\newtheorem{question}[thm]{Question}
%%\newtheorem{prob}[thm]{Problem}
\usepackage{mathtools}
\newcommand{\myabs}[1]{\left|#1\right|}
\newcommand{\mynorm}[1]{\left\|#1\right\|}
\newcommand{\mfrac}[2]{^{#1}/_{#2}}


\title{Simulation of COVID-19 epidemics}
\subtitle{using \texttt{hhh4} model and \texttt{SEIR} model}
\author{Alex Leontiev}
\institute{Datawise Inc.}

\begin{document}
\begin{frame}\titlepage\end{frame}
\begin{frame}{Outline}
  \tableofcontents
\end{frame}
\section{Intro}
\begin{frame}{Goal}
	In this slides I would like to present two ways to model the spread of COVID-19 virus.

	The models were selected based on the following criteria:
	\begin{enumerate}
		\item \textit{spatialness}, that is, model has to account for spread both in space and in time;
		\item model has to {\it use the Trip Distribution data} (i.e. origin-destination matrix) 
			as one of its inputs,
			so we can show off our GPS data;
		\item model has to be \textit{easy to implement}, preferably with R or Python reference source
			code available, or maybe with API available as a part of some software package;
	\end{enumerate}
\end{frame}
\section{Results}
\begin{frame}{Results}
\begin{block}{block}
    block
    \end{block}
    \begin{alertblock}{alert block}
    alert block
    \end{alertblock}
    \begin{exampleblock}{example block}
    example block
    \end{exampleblock}
\end{frame}
\end{document}
